\section{Experimentation and Results}
\label{sec:results}

Six pre-registered hypotheses (H1--H6, H1 comprising two tests, H5 covering the whole of RQ4) are evaluated against the criteria fixed before any result was computed. Reported first is H2, the causal fine-tuning probe and this study's strongest result, since its outcome frames how the remaining, more equivocal results should be read.

\subsection{H2: Fine-Tuning Inflates the Static Benchmark, Not the Dynamic One}
\label{sec:results-h2}
Qwen2.5-Coder-1.5B-Instruct, LoRA-fine-tuned directly on ULT's own repair-verified ground truth (\HTwoSeenN{} problems), gains \HTwoSeenGain{} percentage points (95\% CI \HTwoSeenGainCI{}) on that same seen-ULT split, with \HTwoPreMutkTotal{} of \HTwoPreMutkTotal{} post-fine-tune completions byte-identical to their training target -- direct evidence of memorisation, not transfer. A generic-improvement control, the non-CC-matched dynamic-benchmark tiers (L1--L3) the fine-tuning run never trained on, moves by only \HTwoLOneLThreeGain{} points (CI \HTwoLOneLThreeGainCI{}) over the same fine-tuning run, and the CC-matched dynamic-benchmark control (L4) moves by exactly \HTwoLFourGain{} points. Figure~\ref{fig:onset} shows the full pre/post trajectory across all five populations. The leaked twin PLT, evaluated on the same checkpoint, shows a smaller but still real difference-in-differences gain of \HTwoPltDid{} points (CI \HTwoPltDidCI{}) after subtracting the L1--L3 control, consistent with leaked content providing a partial, indirect memorisation benefit even though the fine-tuning run never saw PLT directly. % cite:external -- "1.5" is the model's own parameter-count identifier (Qwen2.5-Coder-1.5B-Instruct), not a result value

\begin{figure}[t]
\centering
\includegraphics[width=\columnwidth]{comparison_c_onset_curve}
\caption{Pre/post fine-tuning trajectory across seen-ULT, unseen-ULT, PLT, and the two dynamic-benchmark controls.}
\label{fig:onset}
\end{figure}

The approved hypothesis specifies mutation score as the outcome metric; \texttt{baseline\_pass\_rate} is reported here instead, the same estimand under a substituted metric, because pre-fine-tuning mutation score is undefined for \HTwoPreMutkUndefined{} of \HTwoPreMutkTotal{} seen-ULT problems -- the base model almost never produces a suite that clears ULT's baseline-pass gate at all, and mutation score presupposes a baseline-passing suite. \textbf{H2 is supported}: contamination inflates the static benchmark's apparent score by two orders of magnitude more than a generic capability improvement does, while the dynamic benchmark registers a difference indistinguishable from zero.

\subsection{H1 Test 1: Dynamic-L4 vs.\ ULT Shows No Directional Effect}
Restricting the dynamic benchmark to its highest complexity tier (L4, cyclomatic complexity ${\geq}8$) matches ULT's own complexity floor, giving a complexity-controlled comparison across \HOneTOneN{} of the nine canonical models (Qwen2.5-Coder-1.5B-Instruct excluded, a genuine floor effect: zero baseline-passing L4 completions). A paired Wilcoxon signed-rank test gives a rank-biserial effect size of \HOneTOneEffectR{} ($p=\HOneTOneP{}$), favouring the dynamic benchmark in direction but not distinguishable from no effect. Pairing is conducted at model level rather than the approved wording's literal model--problem level, a structural necessity rather than a simplification: the dynamic benchmark and ULT are generated from disjoint sources and share no common problem identity to pair on. \textbf{H1 Test 1 is not supported.} This null result is itself informative rather than merely inconclusive: it is evidence against the alternative explanation that the dynamic benchmark's H2 finding (Section~\ref{sec:results-h2}) simply reflects the dynamic benchmark being a harder task in general -- at matched complexity, models that are not contamination-affected score comparably on both benchmarks, which is consistent with H2's near-zero dynamic-benchmark gain reflecting an absence of contamination rather than a ceiling or difficulty artefact. % cite:external -- "1.5" is the model's own parameter-count identifier, not a result value

\subsection{H1 Test 2: PLT Scores Higher Than ULT, Direction-Consistent but Correction-Dependent}
ULT and its leaked twin PLT are compared directly, CC-band matched, across all nine canonical models. PLT scores higher than ULT on every conditioning tested (effect size \HOneTTwoEffectR{}, raw $p=\HOneTTwoP{}$) -- the contamination-consistent direction, since PLT is deliberately the leaked variant of the same underlying problems. Within RQ1's two-test Holm-Bonferroni family (H1 Test 1 and H1 Test 2 jointly), the raw pre-registered-primary reading does not survive correction ($p=\HOneTTwoP{}\rightarrow\HOneTTwoRawHolmP{}$), but the definition-matched co-primary reading does ($p=\HOneTTwoAdjP{}\rightarrow\HOneTTwoAdjHolmP{}$). \textbf{H1 Test 2 is mixed}: direction agreement across every conditioning is the robust claim; whether it clears significance depends on which conditioning and correction is quoted, and both should be reported together rather than either alone.

\subsection{H3: Mutation-Killing Effectiveness Declines Monotonically With Tier}
Spearman correlation between ordinal tier index and mutation score is negative and significant for all nine canonical models, surviving Holm-Bonferroni correction within both the all-tier family (\HThreeAllTierSig{}) and a ceiling-robustness family excluding the near-ceiling L1 tier (\HThreeLTwoLFourSig{} computable models). The approved hypothesis's stronger, literal clause -- that mean mutation score never increases between adjacent tiers -- holds exactly for seven of the nine models; two (GPT-5 and GPT-5-mini) each show a single upward step at the L1$\rightarrow$L2 transition, both under one percentage point and both already at or above 98\% mutation score at L1, consistent with near-ceiling measurement noise rather than a substantive reversal. \textbf{H3 is supported as a monotone trend}, with the literal non-rebound clause holding for seven of nine models.

\subsection{H5: Score-Level Stability Is Strong, Rank-Level Stability Falls Short}
Two properties are evaluated separately rather than blended into a single verdict. Each model's own aggregate mutation score is highly stable across re-instantiations: standard deviation \HFiveSdRange{} for most canonical models. Pairwise Kendall's $\tau$ between per-instantiation model rankings is markedly weaker: mean $\tau=\HFiveNFourTau{}$ at the four-instantiation headline scope, and $\tau=\HFiveAllPairsTau{}$ across all pairs including a partial fifth instantiation, both below the approved \HFiveTauThreshold{} threshold. Only \HFiveNFourSigCount{} of the six headline-scope pairs is individually significant at raw $p<\SigLevel{}$, and \HFiveNFourHolmCount{} survives Holm correction. This shortfall traces to the evaluated model population's own clustering rather than to instrument unreliability: several canonical models' pooled scores sit within a spread smaller than their own between-instantiation standard deviation, so no ranking instrument could reliably order them from this data alone. Reaching the pre-registered fifth instantiation for the two remaining models (both from one provider) was costed directly against this diagnosis and declined, since the same clustering predicts further data would not close the gap. \textbf{H5's threshold is not met}: score-level reliability is strong and not in question; rank-level stability, as specified, falls short.

\subsection{H6: Construct Validity Against an External Proxy}
Dynamic-benchmark mutation score is correlated against published Pass@1 scores on BigCodeBench, an external, independently maintained code-generation benchmark, and the same correlation is computed for ULT. Under the conditional (pre-registered primary) metric, dynamic-benchmark rank agreement with ULT is $\tau=\HSixCondTau{}$, below the pre-registered \HSixBandLow{}--\HSixBandHigh{} validity-supporting band; under the coverage-adjusted co-primary metric, agreement is $\tau=\HSixAdjTau{}$, above that band. Decomposing the coverage-adjusted figure shows baseline-pass-rate agreement alone already accounts for $\tau=\HSixBpTau{}$ of the total -- most of the apparent agreement reflects which models can produce an executing, passing suite at all, not which suites are the better mutation killers once that is controlled for. \textbf{H6 clause 1 is mixed}: positive under both metrics, but the conditional metric's own agreement is weaker than the pre-registered band anticipated, while the coverage-adjusted metric's agreement is strong enough to raise a redundancy concern by the same pre-registered criterion.

A four-model cross-family extension (adding non-Qwen open-source models to break an all-Qwen scaling confound present in the three-model canonical BigCodeBench comparison) shows the canonical reading's near-perfect correlation was substantially a within-family size artefact: raw conditional agreement collapses from a perfect canonical reading to $\tau=\HSixDynBcbRawTau{}$ at seven models, while the coverage-adjusted reading survives largely intact at $\tau=\HSixDynBcbAdjTau{}$ (Holm-adjusted $p=\HSixDynBcbAdjHolmP{}$), shown in Figure~\ref{fig:collapse}. Two of the seven cross-family points are individually thin-cell by this study's own $n_{\text{defined}}{<}20$ threshold (one model's ULT-side sample: \HSixCodeLlamaUltNDefined{} defined rows; another's dynamic- and ULT-side samples: \HSixMistralDynNDefined{} and \HSixMistralUltNDefined{}), a caveat that travels with the cross-family figures. The coverage-adjusted correlation is itself sensitive to mechanical test-file repair: excluding every repaired row moves the same figure from $\tau=\HSixSanitizerWithTau{}$ to $\tau=\HSixSanitizerWithoutTau{}$, still significant but a smaller number to cite when the estimate's robustness matters.

\begin{figure}[t]
\centering
\includegraphics[width=\columnwidth]{collapse_vs_survival}
\caption{Cross-family extension: raw conditional agreement collapses under added model diversity; coverage-adjusted agreement survives.}
\label{fig:collapse}
\end{figure}

Clause 2, whether the dynamic-proxy correlation is weaker than the static-proxy correlation, is answerable directly from data already computed for clause 1: the static (ULT) correlation with BigCodeBench is $\tau=\HSixUltBcbRawTau{}$ raw and $\tau=\HSixUltBcbAdjTau{}$ coverage-adjusted, both weaker than the dynamic benchmark's own $\HSixDynBcbRawTau{}$ and $\HSixDynBcbAdjTau{}$. \textbf{H6 clause 2 is not supported}: the dynamic-proxy correlation is the stronger one under every statistic and metric definition tested, opposite to the pre-registered prediction. Stated plainly, the dynamic benchmark tracks external code-generation capability at least as well as, and by most measures better than, the static benchmark it was designed as a contamination-resistant alternative to. Figure~\ref{fig:heatmap} summarises rank agreement across all three benchmark pairs and four scoring metrics side by side.

\begin{figure}[t]
\centering
\includegraphics[width=\columnwidth]{ranking_matrix_heatmap}
\caption{Kendall $\tau$ rank agreement by metric and benchmark pair, Holm-Bonferroni significance marked.}
\label{fig:heatmap}
\end{figure}

\subsection{H4: Declared Untestable by Construction}
The approved hypothesis predicts a larger static-versus-dynamic score gap at low complexity tiers (L1--L2) than at high ones (L3--L4). ULT and PLT contain no problems below cyclomatic complexity eight, a filter applied by their own source authors, so no static-benchmark score exists at L1--L2 to form the predicted gap with. \textbf{H4 is dropped as untestable by construction}, a decision taken with supervisor sign-off rather than reshaped post hoc; the reason itself is informative, since a complexity-filtered static benchmark cannot pose the question the hypothesis asks. RQ3 is instead answered narratively by H3's within-dynamic-benchmark tier gradient (Section~\ref{sec:results}.D) together with this structural finding.

\subsection{Additional Disclosures}
Gemini 3.1 Pro Preview, one of the nine canonical models, is a preview-track endpoint carrying deprecation risk beyond that of a generally available model; two of the nine canonical models (both from the same provider) did not reach the pre-registered fifth instantiation used in H5, an asymmetry disclosed at the point H5 is reported rather than left implicit in the model roster. ULT's own published ground truth passes this study's harness at a substantially lower rate than its role as a correctness reference would suggest, a finding independently corroborated by an internal recomputation of the benchmark's own complexity labels and by a direct contradiction between the source paper's stated mutation-score formula and its own published scoring code (Section~\ref{sec:methods}.C); none of this study's own reported metrics are affected, since neither of its two primary definitions matches the contradicted formula, but the finding is reported here as a caveat on ULT's suitability as a ground-truth reference more broadly. Equivalent-mutant validation (Section~\ref{sec:methods}.F) is manual, hand-inspected work not yet complete at the time of writing; the raw, unadjusted mutation score is reported as the default headline figure throughout, and the validated companion score will be reported alongside it once the inspection sample is complete. % cite:external -- "3.1" is the model's own version identifier (Gemini 3.1 Pro Preview), not a result value
